\scanpage{150}
\uline{Corollaire} 1 \unsure{Soient} $(E,q),(E',q')$ formes quad. rég. de $\overline{\mathrm m}$ rg \inserted{$N=2n$ ou $N=2n+1$} sur $S$, \deleted{$F\unsure{\in|\subset}E$} et $H\subset[1,n]$, enfin $D\in G_H(E,q)$, $D'\in G_H(E',q')$. \deleted{Alors si} $D=\{E_1,\ldots,E_{\unsure{r|\nu}}\}$, $D'=\{E'_1,\ldots,E'_r\}$. \deleted{Supposons $S$ affine, et supposons} Posons $Q=E_r^\perp/E_r$, $Q'=E_r^{\prime\perp}/E'_r$, de sorte que $q_\theta$ et $q_{\theta'}$ sont des formes quad. de $=$ rang sur $S$. Si $q_\theta\simeq q_{\theta'}$, \inserted{$\unsure{R|P}_i=E_i/E_{i-1}\simeq R'_i=E'_i/E'_{i-1}$,} et si $S$ est affine, alors \deleted{$q\unsure{\simeq}q'$} $\exists\,u:(E,q)\simeq(E',q')$ tel que $u(D)=D'$

On en conclut

\uline{Corollaire} 2 Il existe loc (ét) un isom $u:(E,q)\longrightarrow(E',q')$ tel que $u(D)=D'$.

Donc,

\uline{Corollaire 3} $\underline G_H(E,q)$ est un espace homogène (ét) sous $\mathrm{Aut}(E,q)=\mathrm O(q)$. En d'autres \uline{termes}, si $D\in G_H(E,q)$, alors
\[
\begin{aligned}
&\mathrm O(q)\longrightarrow\underline G_H\\
&u\longmapsto u(D)
\end{aligned}
\]
est un épim (ét), \deleted{identifiant} définissant un isom
\[
\underline G_H\xleftarrow{\sim}\mathrm O(q)/\unsure{\widetilde{\mathrm O}|\mathrm B}_{D}(q)
\]
où $\widetilde{\mathrm O}_{D}(q)$ est le stabilisateur de $D$ dans $\mathrm O(q)$.
% archive sheet number: 149
